% mazzi: 2   % del 02 solo le pp. 1--13 (la rilevanza economica va col mazzo 03)
\chapter{La concessione e l'appalto di servizi pubblici}

Quando l'ente pubblico affida a un soggetto privato --- nel caso tipico una società sportiva ---
la gestione di un impianto, occorre stabilire se l'affidamento sia una \textbf{concessione} o un
\textbf{appalto di servizi}. La distinzione non è nominalistica: da essa dipendono le modalità
concrete di affidamento, che sono diverse nei due casi.

% ---------------------------------------------------------------------------
\section{Concessione o appalto di servizi?}

\subsection{Il costo del servizio}

\textbf{Criterio del costo}: su chi grava il corrispettivo dell'attività.

\begin{itemize}
  \item \textbf{concessione di pubblici servizi}: il costo del servizio grava sugli
    \textbf{utenti} $\rightarrow$ presuppone sia che il servizio sia reso a \textbf{terzi} (e non
    all'Amministrazione), sia che il corrispettivo sia \textbf{in tutto o in parte} a carico degli
    utenti;
  \item \textbf{appalto di servizi}: spetta all'\textbf{Amministrazione} l'onere di compensare
    l'attività svolta dal privato $\rightarrow$ la prestazione è resa in favore dell'\textbf{ente
    aggiudicatore} e \textbf{a titolo oneroso}.
\end{itemize}

\rubrica{Perché ``in tutto o in parte''}

L'Amministrazione può erogare al gestore un \textbf{contributo di gestione}: il privato non fissa
liberamente le tariffe all'utenza, ma deve applicare le \textbf{tariffe calmierate} approvate
dall'Amministrazione e previste in convenzione. Se il Comune, fedele alla vocazione pubblicistica
del servizio, impone tariffe inferiori a quelle di mercato, il contributo serve a garantire
l'\textbf{equilibrio economico-finanziario} della gestione.

\rubrica{Il caso opposto}

Se l'Amministrazione gestisce direttamente l'impianto e appalta i singoli servizi scorporati
(pulizie, ristorazione, punto vendita), quei servizi sono compensati direttamente da essa e il
privato non ha alcun rapporto con l'utenza: ricorre l'\textbf{appalto}.

\subsection{Il rischio di gestione}

\textbf{Rischio di gestione}: è l'elemento caratterizzante la concessione rispetto all'appalto.

\begin{itemize}
  \item la \textbf{concessione} si configura ogni qual volta l'aggiudicatario assuma il rischio
    della gestione economica del servizio prestato, della sua istituzione e gestione;
  \item il concessionario sostiene una \textbf{parte considerevole} dell'alea economica --- la
    parte sostanziale, non necessariamente la totalità, visto l'eventuale contributo pubblico ---
    che altrimenti farebbe capo all'Amministrazione;
  \item nell'\textbf{appalto} il rischio di gestione resta in capo alla Pubblica Amministrazione.
\end{itemize}

\subsection{Il rapporto diretto con l'utenza}

Dottrina e giurisprudenza sono ormai pacifiche: ricorre la \textbf{concessione di pubblico
servizio}, e non l'appalto, quando a un soggetto esterno alla PA è affidata la fornitura di un
servizio rivolto al \textbf{pubblico degli utenti}, con i quali l'affidatario costituisce
\textbf{direttamente rapporti di tipo giuridico} (Cons.\ Stato sez.\ V n.\ 2294/2002; TAR Emilia
Romagna sez.\ Parma 18.09.1995 n.\ 317; TAR Lombardia sez.\ Milano 4.08.2004 n.\ 3242).

\subsection{I caratteri distintivi}

\begin{table}[htbp]
  \centering \small \renewcommand{\arraystretch}{1.3}
  \begin{tabularx}{\textwidth}{|l|X|X|}
    \hline
    & \textbf{Concessione} & \textbf{Appalto} \\
    \hline
    Natura del titolo & unilaterale, con effetto accrescitivo dei poteri del privato & negoziale:
      il privato esercita i poteri propri dei privati \\
    \hline
    Destinatario diretto & utenza pubblica & Pubblica Amministrazione \\
    \hline
    Rischio di gestione & in capo al privato & in capo alla PA \\
    \hline
    Rapporto & trilatere & contratto a prestazioni corrispettive \\
    \hline
    Costo del servizio & ricade sull'utenza & se lo assume la PA \\
    \hline
  \end{tabularx}
  \caption{Concessione e appalto di servizi: caratteri distintivi.}
  \label{tab:concessione-appalto}
\end{table}

\begin{itemize}
  \item \textbf{effetto accrescitivo dei poteri}: nella concessione l'Amministrazione concede al
    privato di erogare in sua sostituzione un servizio pubblico $\rightarrow$ il privato affianca
    ai poteri propri dell'operatore economico anche poteri di natura pubblicistica;
  \item \textbf{appalto}: contratto tipico del codice civile, sinallagmatico, a prestazioni
    corrispettive;
  \item \textbf{rapporto trilatere} della concessione: Amministrazione proprietaria $\rightarrow$
    privato gestore che si sostituisce a essa nell'erogazione $\rightarrow$ utenza, che si fa
    carico in tutto o in parte del costo del servizio.
\end{itemize}

\subsection{La definizione del codice dei contratti pubblici}

\textbf{Concessione di servizi} (art.\ 3 D.Lgs.\ 50/2016, codice dei contratti pubblici):
``contratto a titolo oneroso stipulato per iscritto in virtù del quale una o più stazioni
appaltanti affidano a uno o più operatori economici la fornitura e la gestione di servizi diversi
dall'esecuzione di lavori \ldots\ riconoscendo a titolo di corrispettivo unicamente il diritto di
gestire i servizi oggetto del contratto o tale diritto accompagnato da un prezzo, con assunzione in
capo al concessionario del \textbf{rischio operativo} legato alla gestione dei servizi''.

\begin{itemize}
  \item il ``prezzo'' che può accompagnare il diritto di gestione è il \textbf{contributo
    pubblico};
  \item \textbf{conclusione}: nella gestione di un impianto sportivo ricorre \textbf{generalmente
    la concessione di servizi}, perché a fronte dell'affidamento il corrispettivo consiste anche o
    solo nella gestione dello stesso, e l'alea è in capo al gestore e non alla PA --- l'incognita
    è la risposta dell'utenza all'offerta sportiva, date le tariffe stabilite
    dall'Amministrazione, e quindi il ritorno economico della gestione.
\end{itemize}

% ---------------------------------------------------------------------------
\section{La natura di servizio pubblico nella giurisprudenza}

\textbf{TAR Lombardia sez.\ III n.\ 5021/2009}: i tratti distintivi del servizio pubblico sono
ravvisabili nel servizio di gestione della piscina comunale.

\begin{itemize}
  \item è attività \textbf{oggettivamente correlata} alla realizzazione di interessi pubblici,
    perché funzionale per caratteristiche intrinseche a consentire a qualunque interessato lo
    svolgimento di attività sportiva;
  \item l'attività sportiva è \textbf{strettamente connessa con la tutela della salute}, che
    l'\textbf{art.\ 32 Cost.} individua quale fondamentale diritto dell'individuo e interesse
    della collettività $\rightarrow$ la natura pubblicistica della gestione si àncora anche a
    questo diritto fondamentale, oltre che al carattere sociale del servizio;
  \item è attività prestata \textbf{direttamente in favore degli utenti} per soddisfare interessi
    di rilevanza generale.
\end{itemize}

\textbf{Corte dei Conti, sez.\ controllo Lombardia n.\ 124/2011}: quel principio di fondo deve
trovare concreta attuazione nell'osservanza, da parte del gestore, di \textbf{livelli minimi di
servizio} per quantità e qualità, nonché nella \textbf{continuità e regolarità} del servizio in
favore della collettività degli utenti. La verifica fattuale va compiutamente assolta dall'ente
locale, per configurare correttamente l'affidamento in termini di concessione di pubblico
servizio.

\begin{itemize}
  \item \textbf{livelli quantitativi}: giorni e orari di apertura da fissare nell'atto
    concessorio;
  \item \textbf{livelli qualitativi}: figure qualificate richieste (direttori di impianto,
    istruttori specializzati), coerentemente con la normativa regionale di riferimento;
  \item \textbf{continuità}: un servizio pubblico non può essere interrotto $\rightarrow$
    l'Amministrazione non può limitarsi a concedere il bene, ma deve imporre e verificare questi
    standard.
\end{itemize}

% ---------------------------------------------------------------------------
\section{La disciplina applicabile all'affidamento}

\subsection{Il quadro previgente: l'art.\ 30 D.Lgs.\ 163/2006}

\textbf{Art.\ 30 comma 1 D.Lgs.\ 163/2006} (abrogato dal D.Lgs.\ 50/2016): ``salvo quanto disposto
nel presente articolo, le disposizioni del codice non si applicano alle concessioni di servizi''.

\begin{itemize}
  \item per questa esclusione espressa il D.Lgs.\ 163/2006 era di fatto un \textbf{codice degli
    appalti} e non un vero codice dei contratti pubblici;
  \item per l'affidamento in concessione di servizi pubblici come quello sportivo si riteneva
    quindi legittimamente prevalente la \textbf{normativa regionale};
  \item in mancanza di questa valevano le indicazioni generali della seconda parte dell'art.\ 30
    stesso.
\end{itemize}

\rubrica{Principi della seconda parte dell'art.\ 30}

La scelta del concessionario doveva avvenire nel rispetto dei principi desumibili dal Trattato e
dei principi generali sui contratti pubblici, in particolare:

\begin{multicols}{2}
\begin{itemize}
  \item trasparenza;
  \item adeguata pubblicità;
  \item non discriminazione;
  \item parità di trattamento;
  \item mutuo riconoscimento;
  \item proporzionalità;
  \item previa gara informale.
\end{itemize}
\end{multicols}

Alla gara informale sono invitati \textbf{almeno cinque concorrenti}, se sussistono in tale numero
soggetti qualificati in relazione all'oggetto della concessione, e con \textbf{predeterminazione
dei criteri selettivi}.

\subsection{La norma quadro: art.\ 90 legge 289/2002, comma 25}

\textbf{Art.\ 90 comma 25 L.\ 289/2002}: ``nei casi in cui l'ente pubblico territoriale non intenda
gestire direttamente gli impianti sportivi, la gestione è affidata \textbf{in via preferenziale} a
società e associazioni sportive dilettantistiche, sulla base di \textbf{convenzioni} che ne
stabiliscono i criteri d'uso e previa determinazione di criteri generali e obiettivi per
l'individuazione dei soggetti affidatari. Le \textbf{Regioni disciplinano, con propria legge, le
modalità di affidamento}''.

\begin{itemize}
  \item a differenza dell'ordinamento sportivo, privo di legge quadro, in materia di impianti
    sportivi esiste dunque una \textbf{vera norma quadro}, ed è quella che ha consentito finora
    alle Regioni di legiferare;
  \item il criterio di \textbf{preferenzialità} è di difficile decodificazione --- le Regioni che
    hanno provato a declinarlo lo hanno fatto con fatica --- ed è stato attenzionato
    dall'\textbf{Unione Europea} come norma che violerebbe la concorrenza e la \textit{par
    condicio} fra operatori economici.
\end{itemize}

Sul piano operativo la preferenzialità non è mai \textbf{affidamento diretto}: resta ferma in ogni
caso la necessità di una \textbf{pubblicità preventiva}. Le Regioni che le hanno dato il contenuto
più esplicito --- Molise, Toscana e Veneto --- la traducono in una \textbf{priorità procedurale}:
una prima procedura riservata ai soggetti sportivi, e l'apertura ad altri soggetti solo in caso di
\textbf{esito infruttuoso}; soluzione a sua volta già censurata dalla Commissione europea per
violazione dei principi sulla concorrenza (cap.~4).

\subsection{Gli orientamenti giurisprudenziali}

\rubrica{Obbligo di evidenza pubblica}

\textbf{Cons.\ Stato sez.\ V 29 dicembre 2009 n.\ 8914} (su L.R.\ Lombardia 27/2006 art.\ 2):
afferma sia la \textbf{necessità di una procedura ad evidenza pubblica} per l'affidamento in
gestione degli impianti sportivi comunali, sia la possibilità per gli interessati di chiedere
l'inserimento nel bando di \textbf{requisiti di selezione ulteriori}, logicamente connessi
all'oggetto dell'appalto, fermi i limiti di logicità, ragionevolezza, pertinenza e congruità
rispetto allo scopo perseguito.

\rubrica{Affidamento diretto e impianto privo di rilevanza economica}

\textbf{Cons.\ Stato sez.\ V 27 settembre 2011 n.\ 5379} (su L.R.\ Lombardia 27/2006 art.\ 5):
stante l'esatta qualificazione dell'impianto come \textbf{privo di rilevanza economica}, il Comune
ha correttamente proceduto ad \textbf{affidamento diretto} ex art.\ 5 comma 2 della legge
regionale, avendo comunque rispettato i principi di \textbf{trasparenza e imparzialità} imposti
dalla disciplina comunitaria.

\begin{itemize}
  \item \textbf{il caso}: una società sportiva non affidataria impugna l'affidamento diretto
    disposto dal Comune, prima davanti al TAR Lombardia e poi in appello;
  \item \textbf{i due elementi decisivi}: la natura di impianto privo di rilevanza economica, e il
    fatto che l'affidataria fosse una \textbf{polisportiva comunale} aperta a tutte le
    associazioni sportive del territorio, espressamente invitate, che ben potevano e possono
    aderirvi in qualunque momento $\rightarrow$ pur essendo l'atto finale un affidamento diretto,
    a monte si era cercato di riunire in un unico soggetto tutti gli operatori interessati.
\end{itemize}

\textbf{La disciplina lombarda su cui le due pronunce si fondano} (cap.~4): la L.R.\ 27/2006 è la
legge regionale che più di ogni altra \textbf{differenzia la procedura di selezione} a seconda che
l'impianto abbia o meno rilevanza economica. Per gli impianti \textbf{con} rilevanza economica che
per dimensioni e altre caratteristiche richiedono una gestione di tipo imprenditoriale, le società
e associazioni sportive dilettantistiche devono dimostrare di possedere i \textbf{requisiti
imprenditoriali e tecnici} necessari; per quelli \textbf{senza} rilevanza economica è ammesso
l'\textbf{affidamento diretto}, anche agli \textbf{utilizzatori} degli impianti stessi.

\rubrica{Rinnovo e proroga della concessione}

\textbf{Cons.\ Stato sez.\ V n.\ 6132/2011}:

\begin{itemize}
  \item il \textbf{rinnovo} di una concessione può essere legittimamente disposto bandendo una
    gara, essendo le PA sempre assoggettate all'obbligo di procedure ad evidenza pubblica per
    individuare il contraente;
  \item il \textbf{precedente concessionario} va posto sullo stesso piano di ogni altro soggetto
    richiedente --- non vanta alcun diritto di insistenza o di prelazione --- e la gara può essere
    indetta senza particolare motivazione sul diniego del rinnovo;
  \item lo \textbf{stesso principio} vale quando il gestore si offra di accollarsi la
    realizzazione di \textbf{interventi strutturali} (copertura o realizzazione della tribuna,
    rifacimento degli spogliatoi) chiedendo in cambio una \textbf{proroga}: l'operazione non è
    ammissibile se la possibilità di proroga non era esplicitamente prevista nella procedura di
    affidamento originaria, così che ogni altro operatore economico potesse conoscere fin dal
    bando la durata potenzialmente più lunga dell'affidamento.
\end{itemize}
